\chapter{حروف‌نگاری}
\label{cha:typography}
\epigraph{حروف‌نگاری هنر آن است که به زبان آدمی شکلی ماندگار بدهیم.}{رابرت برینگهرست}

متن اصلی با قلم \lr{XM Yekan} در اندازه‌ی ده و با فاصله‌ی خطی ۱٫۲۰ چیده شده
است. عنوان‌ها از همان قلم می‌آیند اما ضخامتشان تغییر می‌کند: بخش ضخیم‌تر از
زیربخش و زیربخش ضخیم‌تر از متن. همین کاهش تدریجی است که صفحه را شبیه کتاب
می‌کند، نه شبیه اسلاید.

\section{وزن‌ها و تأکید}

\lr{XM Yekan} تنها یک وزن دارد، پس ضخامت‌های این قالب با
\lr{AutoFakeBold} ساخته می‌شوند. اگر قلمی چندوزنی مانند وزیرمتن را در
\lr{\texttt{fonts/}} بگذارید و \lr{\texttt{styles/fonts.sty}} را به آن
اشاره دهید، وزن‌ها واقعی می‌شوند. متن لاتین داخل جمله با قلم \lr{Inter}
چیده می‌شود؛ برای نمونه \lr{Clean Architecture} یا
\lr{Domain-Driven Design}.

ضخامت‌های نمایشی مستقیماً هم در دسترس‌اند:
{\displaybook معمولی}، {\displaymed متوسط}، {\displaysemi نیمه‌سیاه}،
{\displaybold سیاه}، {\displayblack خیلی سیاه}.

\subsection{عناصر درون‌متنی}

کلید \kbd{Ctrl}+\kbd{Shift}+\kbd{P} را بزنید تا پنجره‌ی فرمان باز شود، سپس
\cmd{git rebase --interactive} را اجرا کنید. این نسخه \pill[green]{پایدار}
است، رابط قدیمی \pill[red]{منسوخ} شده و مسیر جریانی هنوز
\outlinepill[teal]{آزمایشی} است. آنچه هنگام خواندن \hl{نشانه‌گذاری} می‌کنید
نشانه‌گذاری‌شده می‌ماند.

برچسب‌ها مثل اوبسیدین کار می‌کنند: \tagpill{\#جبر-خطی}
\tagpill[purple]{\#ناتمام} \tagpill[teal]{\#مرجع}.

\subsubsection{عنوانی آرام‌تر}

زیرزیربخش‌ها در اندازه‌ی متن و با رنگ کم‌رنگ‌تر چیده می‌شوند. کارشان
سامان‌دادن به یک بحث طولانی است، بی‌آنکه آن را قطع کنند.

\paragraph{عنوان درون‌خطی} برای یادداشت‌های پرتعریف مفید است، جایی که یک
عنوان کامل بیش از ارزش آن ایده جا می‌گیرد.

\section{حاشیه}

حاشیه‌ی بیرونی ۳۴ میلی‌متر است و این پهنا عمدی
است.\sidenote{یادداشت‌های حاشیه اینجا می‌نشینند. آن‌ها را در حد یک یا دو
جمله نگه دارید.} یادداشت‌های کناری را در خود جا می‌دهد و به ستون متن فضای
نفس‌کشیدن می‌دهد.

\divider

پانوشت‌ها ریز و کم‌صدا هستند و خط کوتاهی بالایشان کشیده می‌شود.\footnote{مثل
همین. در اندازه‌ی ۷٫۴ و چسبیده به پای صفحه.}

\begin{quotation}
نقل‌قول برجسته بدون کادر و با رنگ ملایم‌تر می‌آید. تأکیدش را از فضای سفید
اطرافش می‌گیرد، نه از یک قاب.
\end{quotation}

\section{ریاضیات}

فرمول‌ها با \lr{Latin Modern Math} چیده می‌شوند و مثل هر متن لاتین دیگری از
چپ به راست می‌مانند.

\begin{definition}[همگرایی]
دنباله‌ی $(x_n)$ در فضای متریک $(X, d)$ به $x \in X$ همگراست هرگاه برای هر
$\varepsilon > 0$ عددی مانند $N$ باشد که برای هر $n > N$ داشته باشیم
$d(x_n, x) < \varepsilon$.
\end{definition}

\begin{theorem}[یکتایی حد]
هر دنباله در فضای هاوسدورف حداکثر یک حد دارد.
\end{theorem}

\begin{proof}
فرض کنید $x_n \to x$ و $x_n \to y$ با $x \neq y$. دو همسایگی جدا از هم
$U \ni x$ و $V \ni y$ بگیرید. دنباله از جایی به بعد در هر دو می‌افتد، که با
$U \cap V = \emptyset$ ناسازگار است.
\end{proof}

فرمول نمایشی فضای خودش را می‌گیرد:
\[
  \int_{-\infty}^{\infty} e^{-x^2}\,\mathrm{d}x = \sqrt{\pi}.
\]
